% ============================================================================
%  6  总结与进一步问题
% ============================================================================
\section{总结与进一步的问题}

本文围绕「典范丛——切丛——法丛——共轭 Stiefel--Whitney 类」这条线索，为实
Grassmann 流形的欧氏非嵌入问题建立了一条统一的分析链。对
$\GK=G_k(\R^{k+n})$，典范分解 $\cG_k\oplus\cGc_k\cong\eps^{k+n}$ 与切丛公式
\begin{equation*}
  T\GK\cong\Hom(\cG_k,\cGc_k)
\end{equation*}
把几何问题转化为向量丛特征类的计算问题。若 $M^{d}$ 嵌入 $\R^{d+c}$，则法丛满足
$w(\nu)=w(TM)^{-1}$，非零的共轭类 $\wbar_j(TM)$ 至少迫使 $c\ge j$；再调用欧氏
嵌入的临界次数判据（定理 \ref{thm:critical}），可排除 $c=j$ 的端点。这两条都是
嵌入的\emph{必要}条件，而且由于嵌入包含全局单射性，它们不能替代对嵌入本身的
构造。

本文核验并复现了两类文献结果。第一类是 Oproiu 关于 $G_{2,n}$ 与 $G_{3,n}$ 的
非嵌入定理：其二进制参数分段、稳定切丛公式以及关键 Schubert 类的非零性，都必须
在真实的 Grassmann 上同调环中理解；尤其要把临界次数判据与一般的秩截断分开。
第二类是 Hiller--Stong 的浸入维数结果：Lam 型上界、命题 4.1、未编号的
\emph{Corollary}、命题 4.2 与命题 4.3 说明模 $2$ 特征类可以给出一系列 immersion
下界，并须注意它们与 Whitney 型上界并不是同一个数。这两类结论的研究对象不同
——前者是 embedding，后者是 immersion——不能互相替代。

本文的个人复现包括：Grassmann 切丛公式的局部图像证明、法丛秩障碍定理、
$G_{2,2}$ 的商环与 Gr\"obner 标准单项式检查、$G_{2,3}$ 的附加交叉检查，以及
$G_{3,6}$ 边界类的商环秩核验。这些复现的价值在于把「形式展开——商环约化——
非零性判定——维数排除」四个步骤分开检查；它们不构成新的非嵌入定理。凡缺乏一般
证明的固定参数计算，本文只标为文献结果或个人复现，不外推为定理。

归结起来，全文的逻辑层级是清楚的。第一层是向量丛结构：局部图像坐标给出
$T\GK\cong\Hom(\cG_k,\cGc_k)$。第二层是法丛的必要条件：秩截断给 $c\ge j$，
临界次数判据把它加强为 $c\ge j+1$。第三层是具体参数族的计算：Oproiu 用二进制
分段与 Schubert 非零性得到非嵌入结论，Hiller--Stong 用稳定切丛与模 $2$ 计算得到
浸入维数下界。第四层是可复现性：Hiller、Hiller--Stong 与 Stong 的商环、旗流形
方法说明，任何「非零」断言都必须在真实的 Grassmann 上同调商环中完成约化与评价。
本文没有把特征类公式包装成「最小嵌入维数已求出」的强结论；特征类给出必要条件，
文献的 Schubert／旗流形计算把必要条件落实到参数族，而嵌入的构造或更精细的排除
属于另一个问题。

后续工作仍有两项。其一，对第 \ref{sec:classical} 章 $G_{3,n}$ 的边界参数与第 5 章
的代表性计算做独立的商环复核，避免把自由多项式中的形式非零当成上同调类非零。
其二，在明确标注「探索性观察」的前提下，对 $G_{2,n}$ 的共轭类做有限参数实验，
并与 Oproiu 的参数族逐项比较；在完成文献查重与严格证明之前，不宣称新的闭式
定理。

综上，Stiefel--Whitney 类方法适合提供可靠的非嵌入必要条件与浸入维数下界，但
通常不能单独决定最小嵌入维数。后续研究应保持三条边界：在真实上同调商环中计算；
在引用中给出准确的定理编号与参数范围；始终分开书写 embedding、immersion 与
其他受限嵌入问题。
